\section{Running the Adventure}
\label{sec:runningAdventure}

To run this adventure you will need \handbook{} containing all rules and guidelines on how to run \cosmereRPG{}. This book contains only adventure-specific details and suggestions on how to run it.

\quoteBox{
	You can read (or paraphrase) a text in such a box to your players. Alternatively, you can describe the scene using your own words.
}

Whenever a characters name appears in \enemy{bold}, it means there is corresponding stat block available. You can use it to get more info about the character's attributes and abilities.

Items appearing in \loot{bold italic} have special rules you can find at the end of this book.

\reasonBox{
	Boxes like this will help you understand the reason behind some of the design choices. Those boxes have no impact on the game, but they may help you modify some elements of the adventure if you so choose.
}

\ptrBox{
	When a plot might allow for a certain Radiant Spren to be attracted by the player's actions, that will be indicated with such a box. However, use your own judgment whenever a players does something not accounted by this book.
}

\spoilerBox{
	On many occasions, this book will reference events or characters from \bookSA{} series. This box will warn you of spoilers, but also it will provide original source if you wished to read more details about what's referenced from official books.
}

If you are using this book in electronic version, it has also several quality-of-life improvements to make your life easier when you navigate this book:

Whenever a name of a chapter or section appears (e.g. "\nameref{sec:runningAdventure}"), you can click on it to quickly go to the chapter in question.

All characters' names (e.g. \npcOdium{}) are links that lead to \hyperref[personae]{Dramatis Person\ae} chapter, where all the NPCs are listed.

Unique items (e.g. \lootExample{}) and adversaries (e.g. \enemyAxehound{}) are also links that will lead you to the page with their unique rules. 



\subsection{Adjusting Difficulty}

This book assumes your players will form a party of TODO characters starting at TODO level. However, it is perfectly fine to have a different size of your group and/or to start playing at a different level (e.g. because the characters already shared another adventure together). Also, it is possible that challenges presented in this book are not aligning with strong suits of the party.

If you see that the balance is a bit off, feel free to adjust the difficulty. You can change a number or tier of the enemies, modify endeavors' requirements or weave some other changes into the story.



\subsection{Handling Canon}
\label{ssec:runningAdventureCanon}

Despite treating original source with respect, this adventure is not canon. While playing the game your primary focus should be fun. For everyone, "fun" means something different - either following official lore to the letter, creating your own touching story, or just hitting bad guys with your sword. Make sure you know what is fun for you and your group and make it your primary objective. If that means modifying this adventure or straying away from the canon - go for it!

You can find more information on that topic in “Canon and Continuity” section of \handbook{}.




\subsection{Playing with Singers}

TODO; \lipsum[1][1-5]



\subsection{Playing with Radiants}

TODO; \lipsum[1][1-5]


\subsection{Playing with TODO}

TODO; \lipsum[1][1-5]



\subsection{Unique Rules}

TODO; \lipsum[1][1-5]


\subsubsection{\abilityExampleName{}}
\label{ssec:ruleExample}

TODO; \lipsum[1][1-5]

\abilityExample{}

\lipsum[1]


